\section{账本事件的区域制度深描：eastern-jin-sixteen-kingdoms}

\label{sec:ledger-depth-eastern-jin-sixteen-kingdoms}

本组叙事以本卷事件账本的可核年序为骨架，将政治节点放回地方资源、制度中介、人口流动与材料类型中。每条来源能够证明的范围不同，故不以朝廷命令、战报数字或后出传记替代基层社会的完整经验。

\subsection{地方资源与制度接口}

账本条目\texttt{JH-0307-SIMA-RUI-JIANYE}记载，永嘉元年（307年），西晋江东发生“司马睿出镇建业，王导等协助联络江东士族并重整东南军政”。这一节点可放在西晋北方内战使东南成为宗室可依托的相对稳定区域，江东旧族又需防范流民与割据军的条件下理解；其后续影响被账本概括为建业逐渐形成晋室南方中枢，为永嘉之乱后的南渡政权提供制度与人脉基础。就地方尺度而言，事件并不只属于朝廷或统帅：征发、道路、仓储、户籍、市场、寺院与家庭关系都可能决定命令如何抵达、战事如何延长，或接收何以在不同地区呈现不一样的速度。

本条列举的核心材料为《晋书·元帝纪》；《晋书·王导传》；《建康实录》。这些材料较适合钉住事件的发生、官方表述或特定人物、地点，却不自动构成对全境人口、收入、兵数、信仰或动机的调查。账本的置信与材料提示是“出镇建业和王导参与可靠；任命缘由、士族最初支持程度与行政控制范围有争议”；来源限度进一步说明“东晋建国前史、建业城市网络与侨寓士族研究”。因此，不能把条目中的政权名称、行政分类或战报修辞直接转译为固定族群和统一社会意志，也不能将制度宣布写成各地同步实施。

在叙事上，更稳妥的做法是把这一事件视为一条需要地方中介才能运作的链：中央的命令、地方官与精英的选择、运输与劳力的条件、以及普通家庭可能面对的迁徙、赋役或安全风险彼此相连。现有材料能够显示其中若干环节，却保留了大量沉默。承认这种沉默并非放弃解释，而是避免以一条编年记录覆盖不同区域和社群的差异。

账本条目\texttt{JH-0313-ZU-TI-NORTHERN-EXPEDITION}记载，建兴元年（313年），东晋前身与河南地方集团发生“祖逖率部北渡长江，经屯田与结盟经营豫州、河南一带”。这一节点可放在北方沦陷迫使流民与将领自行组织防卫，建业政权也需在江北建立缓冲的条件下理解；其后续影响被账本概括为祖逖的行动一度恢复淮北联系，却依赖个人网络而未形成稳定北伐体制。就地方尺度而言，事件并不只属于朝廷或统帅：征发、道路、仓储、户籍、市场、寺院与家庭关系都可能决定命令如何抵达、战事如何延长，或接收何以在不同地区呈现不一样的速度。

本条列举的核心材料为《晋书·祖逖传》；《资治通鉴·晋纪》；《世说新语》注引《晋阳秋》。这些材料较适合钉住事件的发生、官方表述或特定人物、地点，却不自动构成对全境人口、收入、兵数、信仰或动机的调查。账本的置信与材料提示是“北渡、屯田和北伐经营大体可靠；收复地域、兵粮数量和地方集团归附程度有争议”；来源限度进一步说明“祖逖北伐、屯田与永嘉后河南地方秩序研究”。因此，不能把条目中的政权名称、行政分类或战报修辞直接转译为固定族群和统一社会意志，也不能将制度宣布写成各地同步实施。

在叙事上，更稳妥的做法是把这一事件视为一条需要地方中介才能运作的链：中央的命令、地方官与精英的选择、运输与劳力的条件、以及普通家庭可能面对的迁徙、赋役或安全风险彼此相连。现有材料能够显示其中若干环节，却保留了大量沉默。承认这种沉默并非放弃解释，而是避免以一条编年记录覆盖不同区域和社群的差异。

账本条目\texttt{JH-0315-TUOBA-YILU}记载，建兴三年（315年），代国与西晋残余政权发生“拓跋猗卢受晋封为代王，在北部边地扩展部众与城郭控制”。这一节点可放在晋室为牵制汉赵及其他北方势力，需要借助拓跋部的军事力量的条件下理解；其后续影响被账本概括为代国的名号和资源网络为后来北魏兴起保存了政治与人口基础。就地方尺度而言，事件并不只属于朝廷或统帅：征发、道路、仓储、户籍、市场、寺院与家庭关系都可能决定命令如何抵达、战事如何延长，或接收何以在不同地区呈现不一样的速度。

本条列举的核心材料为《魏书·序纪》；《晋书·愍帝纪》；《资治通鉴·晋纪》。这些材料较适合钉住事件的发生、官方表述或特定人物、地点，却不自动构成对全境人口、收入、兵数、信仰或动机的调查。账本的置信与材料提示是“受封与代王称号可靠；封授具体条件、部众规模和对晋朝实际从属有争议”；来源限度进一步说明“代国早期、晋朝册封与北部草原—农耕边地研究”。因此，不能把条目中的政权名称、行政分类或战报修辞直接转译为固定族群和统一社会意志，也不能将制度宣布写成各地同步实施。

在叙事上，更稳妥的做法是把这一事件视为一条需要地方中介才能运作的链：中央的命令、地方官与精英的选择、运输与劳力的条件、以及普通家庭可能面对的迁徙、赋役或安全风险彼此相连。现有材料能够显示其中若干环节，却保留了大量沉默。承认这种沉默并非放弃解释，而是避免以一条编年记录覆盖不同区域和社群的差异。

账本条目\texttt{JH-0317-SIMA-RUI-KING-JIN}记载，建武元年三月（317年），东晋发生“司马睿在建康即晋王位，建立承继西晋的南方朝廷框架”。这一节点可放在长安失守后晋室必须在江南建立可运作的继承中心，王导等调和侨姓与江东旧族的条件下理解；其后续影响被账本概括为建康成为新的都城，侨置、军镇和南北人口流动成为东晋长期课题。就地方尺度而言，事件并不只属于朝廷或统帅：征发、道路、仓储、户籍、市场、寺院与家庭关系都可能决定命令如何抵达、战事如何延长，或接收何以在不同地区呈现不一样的速度。

本条列举的核心材料为《晋书·元帝纪》；《晋书·王导传》；《建康实录》。这些材料较适合钉住事件的发生、官方表述或特定人物、地点，却不自动构成对全境人口、收入、兵数、信仰或动机的调查。账本的置信与材料提示是“即晋王位和建康朝廷形成可靠；拥立者排序、对愍帝消息的反应与实际控制地域有争议”；来源限度进一步说明“东晋建国、建康朝廷与南渡士族联盟研究”。因此，不能把条目中的政权名称、行政分类或战报修辞直接转译为固定族群和统一社会意志，也不能将制度宣布写成各地同步实施。

在叙事上，更稳妥的做法是把这一事件视为一条需要地方中介才能运作的链：中央的命令、地方官与精英的选择、运输与劳力的条件、以及普通家庭可能面对的迁徙、赋役或安全风险彼此相连。现有材料能够显示其中若干环节，却保留了大量沉默。承认这种沉默并非放弃解释，而是避免以一条编年记录覆盖不同区域和社群的差异。

账本条目\texttt{JH-0318-YUAN-EMPEROR}记载，太兴元年三月（318年），东晋发生“司马睿称帝，是为元帝，东晋以建康为都正式确立”。这一节点可放在西晋帝统已在北方断绝，晋王政权需要以帝号整合流民、州镇和旧官僚的条件下理解；其后续影响被账本概括为东晋确立南朝皇统，但其财政、兵源和北伐能力仍受区域分裂限制。就地方尺度而言，事件并不只属于朝廷或统帅：征发、道路、仓储、户籍、市场、寺院与家庭关系都可能决定命令如何抵达、战事如何延长，或接收何以在不同地区呈现不一样的速度。

本条列举的核心材料为《晋书·元帝纪》；《晋书·王导传》；《资治通鉴·晋纪》；《建康实录》。这些材料较适合钉住事件的发生、官方表述或特定人物、地点，却不自动构成对全境人口、收入、兵数、信仰或动机的调查。账本的置信与材料提示是“称帝、改元和建康定都可靠；受禅礼仪、北方士族代表性和南方地方社会的接受程度有争议”；来源限度进一步说明“东晋开国礼仪、侨姓政治与江南国家形成研究”。因此，不能把条目中的政权名称、行政分类或战报修辞直接转译为固定族群和统一社会意志，也不能将制度宣布写成各地同步实施。

在叙事上，更稳妥的做法是把这一事件视为一条需要地方中介才能运作的链：中央的命令、地方官与精英的选择、运输与劳力的条件、以及普通家庭可能面对的迁徙、赋役或安全风险彼此相连。现有材料能够显示其中若干环节，却保留了大量沉默。承认这种沉默并非放弃解释，而是避免以一条编年记录覆盖不同区域和社群的差异。

账本条目\texttt{JH-0319-SHI-LE-LATER-ZHAO}记载，太兴二年（319年），后赵发生“石勒在襄国称赵王并建立独立政权，后赵由汉赵将领集团分出”。这一节点可放在刘聪死后汉赵内部多变，石勒掌握河北军队和地方资源而不愿受平阳节制的条件下理解；其后续影响被账本概括为北方形成汉赵、后赵及东晋并立格局，河北成为后赵长期基地。就地方尺度而言，事件并不只属于朝廷或统帅：征发、道路、仓储、户籍、市场、寺院与家庭关系都可能决定命令如何抵达、战事如何延长，或接收何以在不同地区呈现不一样的速度。

本条列举的核心材料为《晋书·石勒载记》；《资治通鉴·晋纪》；《十六国春秋》辑佚。这些材料较适合钉住事件的发生、官方表述或特定人物、地点，却不自动构成对全境人口、收入、兵数、信仰或动机的调查。账本的置信与材料提示是“称王、都襄国和与汉赵决裂可靠；部众构成、称号变化与控制范围有争议”；来源限度进一步说明“后赵建国、襄国城市与羯胡军事集团研究”。因此，不能把条目中的政权名称、行政分类或战报修辞直接转译为固定族群和统一社会意志，也不能将制度宣布写成各地同步实施。

在叙事上，更稳妥的做法是把这一事件视为一条需要地方中介才能运作的链：中央的命令、地方官与精英的选择、运输与劳力的条件、以及普通家庭可能面对的迁徙、赋役或安全风险彼此相连。现有材料能够显示其中若干环节，却保留了大量沉默。承认这种沉默并非放弃解释，而是避免以一条编年记录覆盖不同区域和社群的差异。

账本条目\texttt{JH-0319-LIU-YAO-RENAMES-ZHAO}记载，光初二年（319年），前赵发生“刘曜在长安改汉国号为赵，以关中为中心重建政权身份”。这一节点可放在石勒在河北自立后，刘曜需要以关中行政与新国号凝聚原汉赵力量的条件下理解；其后续影响被账本概括为前赵与后赵的对峙使北方西东两大军事中心长期竞争。就地方尺度而言，事件并不只属于朝廷或统帅：征发、道路、仓储、户籍、市场、寺院与家庭关系都可能决定命令如何抵达、战事如何延长，或接收何以在不同地区呈现不一样的速度。

本条列举的核心材料为《晋书·刘曜载记》；《资治通鉴·晋纪》；《十六国春秋》辑佚。这些材料较适合钉住事件的发生、官方表述或特定人物、地点，却不自动构成对全境人口、收入、兵数、信仰或动机的调查。账本的置信与材料提示是“改国号与长安中心地位可靠；改号时间、正统论述和前赵与后赵命名的后世区分有争议”；来源限度进一步说明“前赵国号、关中政治与十六国命名史研究”。因此，不能把条目中的政权名称、行政分类或战报修辞直接转译为固定族群和统一社会意志，也不能将制度宣布写成各地同步实施。

在叙事上，更稳妥的做法是把这一事件视为一条需要地方中介才能运作的链：中央的命令、地方官与精英的选择、运输与劳力的条件、以及普通家庭可能面对的迁徙、赋役或安全风险彼此相连。现有材料能够显示其中若干环节，却保留了大量沉默。承认这种沉默并非放弃解释，而是避免以一条编年记录覆盖不同区域和社群的差异。

账本条目\texttt{JH-0320-DU-TAO-DEFEATED}记载，大兴三年（320年），东晋与荆湘地方集团发生“杜弢领导的荆湘叛乱被王敦等镇压，东晋重新取得长江中游部分军政控制”。这一节点可放在永嘉后人口流动、州郡征发和地方军阀竞争使荆湘秩序恶化的条件下理解；其后续影响被账本概括为王敦的军事声望上升，东晋对长江中游的控制更依赖强藩。就地方尺度而言，事件并不只属于朝廷或统帅：征发、道路、仓储、户籍、市场、寺院与家庭关系都可能决定命令如何抵达、战事如何延长，或接收何以在不同地区呈现不一样的速度。

本条列举的核心材料为《晋书·杜弢传》；《晋书·王敦传》；《资治通鉴·晋纪》。这些材料较适合钉住事件的发生、官方表述或特定人物、地点，却不自动构成对全境人口、收入、兵数、信仰或动机的调查。账本的置信与材料提示是“叛乱与其被平定可靠；流民、蛮獠与地方豪强的参与程度及伤亡有争议”；来源限度进一步说明“东晋初年荆湘叛乱、流民武装与区域治理研究”。因此，不能把条目中的政权名称、行政分类或战报修辞直接转译为固定族群和统一社会意志，也不能将制度宣布写成各地同步实施。

在叙事上，更稳妥的做法是把这一事件视为一条需要地方中介才能运作的链：中央的命令、地方官与精英的选择、运输与劳力的条件、以及普通家庭可能面对的迁徙、赋役或安全风险彼此相连。现有材料能够显示其中若干环节，却保留了大量沉默。承认这种沉默并非放弃解释，而是避免以一条编年记录覆盖不同区域和社群的差异。

账本条目\texttt{JH-0321-WANG-DUN-FIRST-REBELLION}记载，大兴四年（321年），东晋发生“王敦以清君侧为名自武昌东下，迫使元帝调整中枢人事并承认其军事优势”。这一节点可放在东晋初建依赖王氏军政网络，元帝试图培植刘隗、刁协等近臣触发强藩反弹的条件下理解；其后续影响被账本概括为皇权对江上军镇的约束明显削弱，建康政治长期受外镇压力塑造。就地方尺度而言，事件并不只属于朝廷或统帅：征发、道路、仓储、户籍、市场、寺院与家庭关系都可能决定命令如何抵达、战事如何延长，或接收何以在不同地区呈现不一样的速度。

本条列举的核心材料为《晋书·元帝纪》；《晋书·王敦传》；《晋书·王导传》；《资治通鉴·晋纪》。这些材料较适合钉住事件的发生、官方表述或特定人物、地点，却不自动构成对全境人口、收入、兵数、信仰或动机的调查。账本的置信与材料提示是“王敦起兵和朝廷让步可靠；王敦具体诉求、王导调停效果和各州响应有争议”；来源限度进一步说明“王敦之乱、荆州军镇与东晋君臣共治研究”。因此，不能把条目中的政权名称、行政分类或战报修辞直接转译为固定族群和统一社会意志，也不能将制度宣布写成各地同步实施。

在叙事上，更稳妥的做法是把这一事件视为一条需要地方中介才能运作的链：中央的命令、地方官与精英的选择、运输与劳力的条件、以及普通家庭可能面对的迁徙、赋役或安全风险彼此相连。现有材料能够显示其中若干环节，却保留了大量沉默。承认这种沉默并非放弃解释，而是避免以一条编年记录覆盖不同区域和社群的差异。

账本条目\texttt{JH-0322-WANG-DUN-TAKES-JIANKANG}记载，永昌元年（322年），东晋发生“王敦第二次起兵攻入建康，杀逐元帝近臣并控制朝政”。这一节点可放在元帝试图重建禁军和限制王敦，引发荆州军政集团以武力进入都城的条件下理解；其后续影响被账本概括为东晋皇权遭重创，王敦虽未废帝却成为决定朝局的强人。就地方尺度而言，事件并不只属于朝廷或统帅：征发、道路、仓储、户籍、市场、寺院与家庭关系都可能决定命令如何抵达、战事如何延长，或接收何以在不同地区呈现不一样的速度。

本条列举的核心材料为《晋书·元帝纪》；《晋书·王敦传》；《晋书·周顗传》；《资治通鉴·晋纪》。这些材料较适合钉住事件的发生、官方表述或特定人物、地点，却不自动构成对全境人口、收入、兵数、信仰或动机的调查。账本的置信与材料提示是“王敦入建康及中枢清洗可靠；城内战斗、元帝应对和死者名单有争议”；来源限度进一步说明“建康政变、王敦集团与东晋中枢政治研究”。因此，不能把条目中的政权名称、行政分类或战报修辞直接转译为固定族群和统一社会意志，也不能将制度宣布写成各地同步实施。

在叙事上，更稳妥的做法是把这一事件视为一条需要地方中介才能运作的链：中央的命令、地方官与精英的选择、运输与劳力的条件、以及普通家庭可能面对的迁徙、赋役或安全风险彼此相连。现有材料能够显示其中若干环节，却保留了大量沉默。承认这种沉默并非放弃解释，而是避免以一条编年记录覆盖不同区域和社群的差异。

账本条目\texttt{JH-0323-YUAN-DEATH}记载，永昌元年闰十一月（323年），东晋发生“元帝去世，明帝司马绍即位，王敦集团与建康朝廷的对抗延续”。这一节点可放在王敦控制朝政后元帝政治空间有限，继承发生在强藩威胁未除之际的条件下理解；其后续影响被账本概括为明帝需要依靠地方将领与王氏分化来恢复都城控制。就地方尺度而言，事件并不只属于朝廷或统帅：征发、道路、仓储、户籍、市场、寺院与家庭关系都可能决定命令如何抵达、战事如何延长，或接收何以在不同地区呈现不一样的速度。

本条列举的核心材料为《晋书·元帝纪》；《晋书·明帝纪》；《晋书·王敦传》。这些材料较适合钉住事件的发生、官方表述或特定人物、地点，却不自动构成对全境人口、收入、兵数、信仰或动机的调查。账本的置信与材料提示是“卒年与继位可靠；遗诏、明帝掌握军权的程度和王氏内部立场有争议”；来源限度进一步说明“东晋皇位继承、王敦危机与建康军政研究”。因此，不能把条目中的政权名称、行政分类或战报修辞直接转译为固定族群和统一社会意志，也不能将制度宣布写成各地同步实施。

在叙事上，更稳妥的做法是把这一事件视为一条需要地方中介才能运作的链：中央的命令、地方官与精英的选择、运输与劳力的条件、以及普通家庭可能面对的迁徙、赋役或安全风险彼此相连。现有材料能够显示其中若干环节，却保留了大量沉默。承认这种沉默并非放弃解释，而是避免以一条编年记录覆盖不同区域和社群的差异。

\subsection{道路、边地与人口移动}

账本条目\texttt{JH-0324-WANG-DUN-COLLAPSE}记载，太宁二年（324年），东晋发生“王敦再度举兵但病死于途中，其部众在建康内外被明帝集团击败”。这一节点可放在王敦试图以武力彻底取代建康中枢，但内部将领对僭越前景并不一致的条件下理解；其后续影响被账本概括为东晋暂时收回都城控制，却未能消除大族和军镇对政治的结构性影响。就地方尺度而言，事件并不只属于朝廷或统帅：征发、道路、仓储、户籍、市场、寺院与家庭关系都可能决定命令如何抵达、战事如何延长，或接收何以在不同地区呈现不一样的速度。

本条列举的核心材料为《晋书·明帝纪》；《晋书·王敦传》；《晋书·温峤传》；《资治通鉴·晋纪》。这些材料较适合钉住事件的发生、官方表述或特定人物、地点，却不自动构成对全境人口、收入、兵数、信仰或动机的调查。账本的置信与材料提示是“王敦病死和叛军失败可靠；病因、部将倒戈次序与朝廷清算范围有争议”；来源限度进一步说明“王敦之乱结局、军镇分化与东晋皇权修复研究”。因此，不能把条目中的政权名称、行政分类或战报修辞直接转译为固定族群和统一社会意志，也不能将制度宣布写成各地同步实施。

在叙事上，更稳妥的做法是把这一事件视为一条需要地方中介才能运作的链：中央的命令、地方官与精英的选择、运输与劳力的条件、以及普通家庭可能面对的迁徙、赋役或安全风险彼此相连。现有材料能够显示其中若干环节，却保留了大量沉默。承认这种沉默并非放弃解释，而是避免以一条编年记录覆盖不同区域和社群的差异。

账本条目\texttt{JH-0325-ZU-TI-DEATH}记载，太宁三年（325年），东晋与河南地方集团发生“祖逖去世，其北伐军和河南经营失去核心协调者”。这一节点可放在祖逖的北方据点高度依赖个人威望与地方结盟，建康难以持续供给的条件下理解；其后续影响被账本概括为东晋对黄河以北的影响进一步萎缩，后赵得以扩张河南。就地方尺度而言，事件并不只属于朝廷或统帅：征发、道路、仓储、户籍、市场、寺院与家庭关系都可能决定命令如何抵达、战事如何延长，或接收何以在不同地区呈现不一样的速度。

本条列举的核心材料为《晋书·祖逖传》；《晋书·郗鉴传》；《资治通鉴·晋纪》。这些材料较适合钉住事件的发生、官方表述或特定人物、地点，却不自动构成对全境人口、收入、兵数、信仰或动机的调查。账本的置信与材料提示是“卒年与北伐经营受挫可靠；部众去向、继任安排和北方地方势力的反应有争议”；来源限度进一步说明“祖逖身后北伐网络、屯田军与东晋边防研究”。因此，不能把条目中的政权名称、行政分类或战报修辞直接转译为固定族群和统一社会意志，也不能将制度宣布写成各地同步实施。

在叙事上，更稳妥的做法是把这一事件视为一条需要地方中介才能运作的链：中央的命令、地方官与精英的选择、运输与劳力的条件、以及普通家庭可能面对的迁徙、赋役或安全风险彼此相连。现有材料能够显示其中若干环节，却保留了大量沉默。承认这种沉默并非放弃解释，而是避免以一条编年记录覆盖不同区域和社群的差异。

账本条目\texttt{JH-0327-SU-JUN-REBELLION}记载，咸和二年（327年），东晋发生“历阳将苏峻以拒绝征召为由起兵，联合祖约威胁建康”。这一节点可放在庾亮试图削弱掌兵的流民将领，苏峻担忧被解除兵权而先发制人的条件下理解；其后续影响被账本概括为建康再度面临军事危机，幼帝与庾氏外戚的统治基础被动摇。就地方尺度而言，事件并不只属于朝廷或统帅：征发、道路、仓储、户籍、市场、寺院与家庭关系都可能决定命令如何抵达、战事如何延长，或接收何以在不同地区呈现不一样的速度。

本条列举的核心材料为《晋书·成帝纪》；《晋书·苏峻传》；《晋书·庾亮传》；《资治通鉴·晋纪》。这些材料较适合钉住事件的发生、官方表述或特定人物、地点，却不自动构成对全境人口、收入、兵数、信仰或动机的调查。账本的置信与材料提示是“苏峻、祖约起兵可靠；征召缘由、兵力数字和士族对叛军的态度有争议”；来源限度进一步说明“苏峻之乱、流民帅与东晋外戚政治研究”。因此，不能把条目中的政权名称、行政分类或战报修辞直接转译为固定族群和统一社会意志，也不能将制度宣布写成各地同步实施。

在叙事上，更稳妥的做法是把这一事件视为一条需要地方中介才能运作的链：中央的命令、地方官与精英的选择、运输与劳力的条件、以及普通家庭可能面对的迁徙、赋役或安全风险彼此相连。现有材料能够显示其中若干环节，却保留了大量沉默。承认这种沉默并非放弃解释，而是避免以一条编年记录覆盖不同区域和社群的差异。

账本条目\texttt{JH-0328-SU-JUN-TAKES-JIANKANG}记载，咸和三年（328年），东晋发生“苏峻军攻入建康，劫持成帝并控制宫城和朝廷”。这一节点可放在朝廷军队缺乏协调，苏峻凭历阳水军沿江迅速推进的条件下理解；其后续影响被账本概括为东晋中枢短暂瘫痪，各地军镇以勤王名义重新组织政治联盟。就地方尺度而言，事件并不只属于朝廷或统帅：征发、道路、仓储、户籍、市场、寺院与家庭关系都可能决定命令如何抵达、战事如何延长，或接收何以在不同地区呈现不一样的速度。

本条列举的核心材料为《晋书·成帝纪》；《晋书·苏峻传》；《晋书·温峤传》；《建康实录》。这些材料较适合钉住事件的发生、官方表述或特定人物、地点，却不自动构成对全境人口、收入、兵数、信仰或动机的调查。账本的置信与材料提示是“陷都、劫帝和宫城控制可靠；城战细节、掠夺程度与成帝行动空间有争议”；来源限度进一步说明“苏峻之乱中的都城破坏、幼帝与流民军研究”。因此，不能把条目中的政权名称、行政分类或战报修辞直接转译为固定族群和统一社会意志，也不能将制度宣布写成各地同步实施。

在叙事上，更稳妥的做法是把这一事件视为一条需要地方中介才能运作的链：中央的命令、地方官与精英的选择、运输与劳力的条件、以及普通家庭可能面对的迁徙、赋役或安全风险彼此相连。现有材料能够显示其中若干环节，却保留了大量沉默。承认这种沉默并非放弃解释，而是避免以一条编年记录覆盖不同区域和社群的差异。

账本条目\texttt{JH-0328-SHI-LE-CAPTURES-LIU-YAO}记载，太和三年（328年），后赵与前赵发生“石勒在洛阳附近击败并俘获前赵皇帝刘曜，前赵主力由此崩解”。这一节点可放在前赵与后赵长期争夺河南，刘曜试图以洛阳为支点压制石勒的条件下理解；其后续影响被账本概括为后赵夺得河南与前赵皇帝，北方权力天平迅速向襄国倾斜。就地方尺度而言，事件并不只属于朝廷或统帅：征发、道路、仓储、户籍、市场、寺院与家庭关系都可能决定命令如何抵达、战事如何延长，或接收何以在不同地区呈现不一样的速度。

本条列举的核心材料为《晋书·石勒载记》；《晋书·刘曜载记》；《资治通鉴·晋纪》。这些材料较适合钉住事件的发生、官方表述或特定人物、地点，却不自动构成对全境人口、收入、兵数、信仰或动机的调查。账本的置信与材料提示是“刘曜被俘和战役结局可靠；战场位置、双方兵数及刘曜战败原因有争议”；来源限度进一步说明“前赵后赵战争、洛阳战略与皇帝被俘研究”。因此，不能把条目中的政权名称、行政分类或战报修辞直接转译为固定族群和统一社会意志，也不能将制度宣布写成各地同步实施。

在叙事上，更稳妥的做法是把这一事件视为一条需要地方中介才能运作的链：中央的命令、地方官与精英的选择、运输与劳力的条件、以及普通家庭可能面对的迁徙、赋役或安全风险彼此相连。现有材料能够显示其中若干环节，却保留了大量沉默。承认这种沉默并非放弃解释，而是避免以一条编年记录覆盖不同区域和社群的差异。

账本条目\texttt{JH-0329-LATER-ZHAO-DESTROYS-FORMER-ZHAO}记载，太和四年（329年），后赵与前赵发生“石勒攻取长安并处置刘曜之子刘熙，前赵政权灭亡”。这一节点可放在刘曜被俘后前赵指挥系统瓦解，关中难以获得有效援军的条件下理解；其后续影响被账本概括为后赵势力进入关中，石勒成为北方最强的军事统治者。就地方尺度而言，事件并不只属于朝廷或统帅：征发、道路、仓储、户籍、市场、寺院与家庭关系都可能决定命令如何抵达、战事如何延长，或接收何以在不同地区呈现不一样的速度。

本条列举的核心材料为《晋书·石勒载记》；《晋书·刘曜载记》；《资治通鉴·晋纪》。这些材料较适合钉住事件的发生、官方表述或特定人物、地点，却不自动构成对全境人口、收入、兵数、信仰或动机的调查。账本的置信与材料提示是“前赵灭亡和长安失守可靠；城内抵抗、处置过程和关中地方反应有争议”；来源限度进一步说明“后赵统一华北前期、长安易手与关中战后秩序研究”。因此，不能把条目中的政权名称、行政分类或战报修辞直接转译为固定族群和统一社会意志，也不能将制度宣布写成各地同步实施。

在叙事上，更稳妥的做法是把这一事件视为一条需要地方中介才能运作的链：中央的命令、地方官与精英的选择、运输与劳力的条件、以及普通家庭可能面对的迁徙、赋役或安全风险彼此相连。现有材料能够显示其中若干环节，却保留了大量沉默。承认这种沉默并非放弃解释，而是避免以一条编年记录覆盖不同区域和社群的差异。

账本条目\texttt{JH-0330-SHI-LE-EMPEROR}记载，建平元年（330年），后赵发生“石勒在襄国称帝，后赵以河北为核心建立更完整的官僚与军政体系”。这一节点可放在消灭前赵后石勒拥有河北、河南和关中部分资源，需要以帝号整合文武集团的条件下理解；其后续影响被账本概括为后赵成为北方主要政权，其征发与城镇建设深刻影响黄河流域社会。就地方尺度而言，事件并不只属于朝廷或统帅：征发、道路、仓储、户籍、市场、寺院与家庭关系都可能决定命令如何抵达、战事如何延长，或接收何以在不同地区呈现不一样的速度。

本条列举的核心材料为《晋书·石勒载记》；《十六国春秋》辑佚；《资治通鉴·晋纪》。这些材料较适合钉住事件的发生、官方表述或特定人物、地点，却不自动构成对全境人口、收入、兵数、信仰或动机的调查。账本的置信与材料提示是“称帝与建平改元可靠；官制仿效程度、地方执行和羯胡—汉人官僚关系有争议”；来源限度进一步说明“后赵帝制、襄国都城与多族群统治研究”。因此，不能把条目中的政权名称、行政分类或战报修辞直接转译为固定族群和统一社会意志，也不能将制度宣布写成各地同步实施。

在叙事上，更稳妥的做法是把这一事件视为一条需要地方中介才能运作的链：中央的命令、地方官与精英的选择、运输与劳力的条件、以及普通家庭可能面对的迁徙、赋役或安全风险彼此相连。现有材料能够显示其中若干环节，却保留了大量沉默。承认这种沉默并非放弃解释，而是避免以一条编年记录覆盖不同区域和社群的差异。

账本条目\texttt{JH-0333-SHI-LE-DEATH}记载，建平四年（333年），后赵发生“石勒去世，太子石弘继位，石虎掌握军政实权”。这一节点可放在石勒建立的统治依赖个人威望与军事集团，继承人缺乏独立兵权的条件下理解；其后续影响被账本概括为石虎的夺权使后赵进入高压扩张与内部猜忌并行的阶段。就地方尺度而言，事件并不只属于朝廷或统帅：征发、道路、仓储、户籍、市场、寺院与家庭关系都可能决定命令如何抵达、战事如何延长，或接收何以在不同地区呈现不一样的速度。

本条列举的核心材料为《晋书·石勒载记》；《晋书·石虎载记》；《资治通鉴·晋纪》。这些材料较适合钉住事件的发生、官方表述或特定人物、地点，却不自动构成对全境人口、收入、兵数、信仰或动机的调查。账本的置信与材料提示是“石勒卒、石弘继位与石虎专权可靠；遗诏内容、石弘自主空间和宗室军权分配有争议”；来源限度进一步说明“后赵继承、石虎军权与羯胡政权稳定性研究”。因此，不能把条目中的政权名称、行政分类或战报修辞直接转译为固定族群和统一社会意志，也不能将制度宣布写成各地同步实施。

在叙事上，更稳妥的做法是把这一事件视为一条需要地方中介才能运作的链：中央的命令、地方官与精英的选择、运输与劳力的条件、以及普通家庭可能面对的迁徙、赋役或安全风险彼此相连。现有材料能够显示其中若干环节，却保留了大量沉默。承认这种沉默并非放弃解释，而是避免以一条编年记录覆盖不同区域和社群的差异。

账本条目\texttt{JH-0334-SHI-HU-COUP}记载，建武元年（334年），后赵发生“石虎废杀石弘，自立为天王，后赵皇位转入其支系”。这一节点可放在石虎已控制禁军与外军，石弘无法制衡其叔父的军事网络的条件下理解；其后续影响被账本概括为后赵政治更趋个人化和暴力化，继承危机被暂时压制而非解决。就地方尺度而言，事件并不只属于朝廷或统帅：征发、道路、仓储、户籍、市场、寺院与家庭关系都可能决定命令如何抵达、战事如何延长，或接收何以在不同地区呈现不一样的速度。

本条列举的核心材料为《晋书·石虎载记》；《晋书·石勒载记下》；《资治通鉴·晋纪》。这些材料较适合钉住事件的发生、官方表述或特定人物、地点，却不自动构成对全境人口、收入、兵数、信仰或动机的调查。账本的置信与材料提示是“废立与石虎掌权可靠；石弘死亡方式、宗室支持和称号礼仪有争议”；来源限度进一步说明“后赵宫廷政变、宗室清洗与军人统治研究”。因此，不能把条目中的政权名称、行政分类或战报修辞直接转译为固定族群和统一社会意志，也不能将制度宣布写成各地同步实施。

在叙事上，更稳妥的做法是把这一事件视为一条需要地方中介才能运作的链：中央的命令、地方官与精英的选择、运输与劳力的条件、以及普通家庭可能面对的迁徙、赋役或安全风险彼此相连。现有材料能够显示其中若干环节，却保留了大量沉默。承认这种沉默并非放弃解释，而是避免以一条编年记录覆盖不同区域和社群的差异。

账本条目\texttt{JH-0334-LI-XIONG-DEATH}记载，玉衡二十四年（334年），成汉发生“李雄去世，李班继位后旋即被李期夺权，成汉内部继承冲突公开化”。这一节点可放在李雄长期统治下未能建立无争议的继承秩序，宗室各支均掌有军政资源的条件下理解；其后续影响被账本概括为成汉从相对稳定转入频繁宫廷冲突，削弱其抵御东晋北伐的能力。就地方尺度而言，事件并不只属于朝廷或统帅：征发、道路、仓储、户籍、市场、寺院与家庭关系都可能决定命令如何抵达、战事如何延长，或接收何以在不同地区呈现不一样的速度。

本条列举的核心材料为《晋书·李雄载记》；《晋书·李期载记》；《华阳国志·李特雄期寿势志》。这些材料较适合钉住事件的发生、官方表述或特定人物、地点，却不自动构成对全境人口、收入、兵数、信仰或动机的调查。账本的置信与材料提示是“李雄卒、李班继位与李期夺权可靠；遗命、宗室派系和成都官僚的作用有争议”；来源限度进一步说明“成汉继承、巴蜀宗室与流民政权转型研究”。因此，不能把条目中的政权名称、行政分类或战报修辞直接转译为固定族群和统一社会意志，也不能将制度宣布写成各地同步实施。

在叙事上，更稳妥的做法是把这一事件视为一条需要地方中介才能运作的链：中央的命令、地方官与精英的选择、运输与劳力的条件、以及普通家庭可能面对的迁徙、赋役或安全风险彼此相连。现有材料能够显示其中若干环节，却保留了大量沉默。承认这种沉默并非放弃解释，而是避免以一条编年记录覆盖不同区域和社群的差异。

账本条目\texttt{JH-0335-MURONG-HUANG-YAN}记载，咸康元年（335年），前燕发生“慕容皝在辽西称燕王，前燕政权以棘城为中心扩展”。这一节点可放在慕容部在辽西积累军事实力，后赵的北方统治也难以有效深入东北的条件下理解；其后续影响被账本概括为前燕成为东北重要政权，参与后赵、东晋与高句丽之间的区域竞争。就地方尺度而言，事件并不只属于朝廷或统帅：征发、道路、仓储、户籍、市场、寺院与家庭关系都可能决定命令如何抵达、战事如何延长，或接收何以在不同地区呈现不一样的速度。

本条列举的核心材料为《晋书·慕容皝载记》；《资治通鉴·晋纪》；《十六国春秋》辑佚。这些材料较适合钉住事件的发生、官方表述或特定人物、地点，却不自动构成对全境人口、收入、兵数、信仰或动机的调查。账本的置信与材料提示是“称燕王与前燕政权形成大体可靠；受晋册封关系、部族联盟和疆域范围有争议”；来源限度进一步说明“前燕建国、慕容部联盟与辽西城镇研究”。因此，不能把条目中的政权名称、行政分类或战报修辞直接转译为固定族群和统一社会意志，也不能将制度宣布写成各地同步实施。

在叙事上，更稳妥的做法是把这一事件视为一条需要地方中介才能运作的链：中央的命令、地方官与精英的选择、运输与劳力的条件、以及普通家庭可能面对的迁徙、赋役或安全风险彼此相连。现有材料能够显示其中若干环节，却保留了大量沉默。承认这种沉默并非放弃解释，而是避免以一条编年记录覆盖不同区域和社群的差异。

\subsection{社会关系与行政分类}

账本条目\texttt{JH-0336-MURONG-HUANG-DEFEATS-DUAN}记载，咸康二年（336年），前燕与段部发生“慕容皝击败段辽并整合辽西部众，前燕的军事与人口资源扩大”。这一节点可放在辽西诸部长期竞争牧地、城郭和对后赵的关系，慕容皝寻求消除侧翼威胁的条件下理解；其后续影响被账本概括为前燕强化对辽西通道的掌控，为后续与后赵和高句丽作战提供兵源。就地方尺度而言，事件并不只属于朝廷或统帅：征发、道路、仓储、户籍、市场、寺院与家庭关系都可能决定命令如何抵达、战事如何延长，或接收何以在不同地区呈现不一样的速度。

本条列举的核心材料为《晋书·慕容皝载记》；《资治通鉴·晋纪》；《十六国春秋》辑佚。这些材料较适合钉住事件的发生、官方表述或特定人物、地点，却不自动构成对全境人口、收入、兵数、信仰或动机的调查。账本的置信与材料提示是“战胜段部和吸纳部众可靠；段部降附规模、部族身份和迁徙路线有争议”；来源限度进一步说明“前燕扩张、段部政治与辽西人口重组研究”。因此，不能把条目中的政权名称、行政分类或战报修辞直接转译为固定族群和统一社会意志，也不能将制度宣布写成各地同步实施。

在叙事上，更稳妥的做法是把这一事件视为一条需要地方中介才能运作的链：中央的命令、地方官与精英的选择、运输与劳力的条件、以及普通家庭可能面对的迁徙、赋役或安全风险彼此相连。现有材料能够显示其中若干环节，却保留了大量沉默。承认这种沉默并非放弃解释，而是避免以一条编年记录覆盖不同区域和社群的差异。

账本条目\texttt{JH-0338-SHI-HU-ATTACKS-YAN}记载，咸康四年（338年），后赵与前燕发生“石虎大军进攻前燕，慕容皝依托棘城及周边地形击退来军”。这一节点可放在后赵试图遏制慕容部扩张，前燕则需保住辽西的政治中心的条件下理解；其后续影响被账本概括为后赵未能消灭前燕，东北独立政权得以继续发展。就地方尺度而言，事件并不只属于朝廷或统帅：征发、道路、仓储、户籍、市场、寺院与家庭关系都可能决定命令如何抵达、战事如何延长，或接收何以在不同地区呈现不一样的速度。

本条列举的核心材料为《晋书·石虎载记》；《晋书·慕容皝载记》；《资治通鉴·晋纪》。这些材料较适合钉住事件的发生、官方表述或特定人物、地点，却不自动构成对全境人口、收入、兵数、信仰或动机的调查。账本的置信与材料提示是“后赵进攻与前燕守住核心地区可靠；兵额、战场细节和战损数字有争议”；来源限度进一步说明“后赵—前燕战争、棘城防御与辽西地缘研究”。因此，不能把条目中的政权名称、行政分类或战报修辞直接转译为固定族群和统一社会意志，也不能将制度宣布写成各地同步实施。

在叙事上，更稳妥的做法是把这一事件视为一条需要地方中介才能运作的链：中央的命令、地方官与精英的选择、运输与劳力的条件、以及普通家庭可能面对的迁徙、赋役或安全风险彼此相连。现有材料能够显示其中若干环节，却保留了大量沉默。承认这种沉默并非放弃解释，而是避免以一条编年记录覆盖不同区域和社群的差异。

账本条目\texttt{JH-0338-LI-SHOU-USURPS-CHENGHAN}记载，咸康四年（338年），成汉发生“李寿废李期自立并改国号为汉，成汉政权的宗室结构再次重组”。这一节点可放在李雄死后继承冲突未止，李寿凭军权和宗室身份夺取成都的条件下理解；其后续影响被账本概括为政权虽暂获统一但统治正当性受损，为东晋西征创造机会。就地方尺度而言，事件并不只属于朝廷或统帅：征发、道路、仓储、户籍、市场、寺院与家庭关系都可能决定命令如何抵达、战事如何延长，或接收何以在不同地区呈现不一样的速度。

本条列举的核心材料为《晋书·李寿载记》；《华阳国志·李特雄期寿势志》；《资治通鉴·晋纪》。这些材料较适合钉住事件的发生、官方表述或特定人物、地点，却不自动构成对全境人口、收入、兵数、信仰或动机的调查。账本的置信与材料提示是“废立、李寿称帝和国号变化可靠；李期罪状、成都官僚态度和改号的实际使用有争议”；来源限度进一步说明“成汉晚期政变、国号与巴蜀政治合法性研究”。因此，不能把条目中的政权名称、行政分类或战报修辞直接转译为固定族群和统一社会意志，也不能将制度宣布写成各地同步实施。

在叙事上，更稳妥的做法是把这一事件视为一条需要地方中介才能运作的链：中央的命令、地方官与精英的选择、运输与劳力的条件、以及普通家庭可能面对的迁徙、赋役或安全风险彼此相连。现有材料能够显示其中若干环节，却保留了大量沉默。承认这种沉默并非放弃解释，而是避免以一条编年记录覆盖不同区域和社群的差异。

账本条目\texttt{JH-0341-WANG-XINGZHI-EPITAPH}记载，咸康七年（341年），东晋江南发生“王兴之夫妇墓志下葬，为建康周边侨姓官员家庭、婚姻与官职提供有纪年实物材料”。这一节点可放在南渡士族在建康周边定居并通过官职、婚姻和墓葬维系社会身份的条件下理解；其后续影响被账本概括为出土墓志为仅依正史重建东晋精英网络提供独立的日期和地域锚点。就地方尺度而言，事件并不只属于朝廷或统帅：征发、道路、仓储、户籍、市场、寺院与家庭关系都可能决定命令如何抵达、战事如何延长，或接收何以在不同地区呈现不一样的速度。

本条列举的核心材料为《晋书·成帝纪》；〈王兴之夫妇墓志〉；南京象山王兴之夫妇墓发掘报告。这些材料较适合钉住事件的发生、官方表述或特定人物、地点，却不自动构成对全境人口、收入、兵数、信仰或动机的调查。账本的置信与材料提示是“墓志纪年与墓主姓名可靠；官衔追赠、家世叙述和墓志能否代表更广泛社会有争议”；来源限度进一步说明“东晋墓志、建康侨姓家族与六朝葬制研究”。因此，不能把条目中的政权名称、行政分类或战报修辞直接转译为固定族群和统一社会意志，也不能将制度宣布写成各地同步实施。

在叙事上，更稳妥的做法是把这一事件视为一条需要地方中介才能运作的链：中央的命令、地方官与精英的选择、运输与劳力的条件、以及普通家庭可能面对的迁徙、赋役或安全风险彼此相连。现有材料能够显示其中若干环节，却保留了大量沉默。承认这种沉默并非放弃解释，而是避免以一条编年记录覆盖不同区域和社群的差异。

账本条目\texttt{JH-0342-MURONG-HUANG-GOGURYEO}记载，咸康八年（342年），前燕与高句丽发生“慕容皝进攻高句丽都城丸都，掳掠人口并改变辽东力量关系”。这一节点可放在前燕为控制辽东通道并消除后方威胁，选择对高句丽核心地区发动远征的条件下理解；其后续影响被账本概括为高句丽短期受挫而前燕获得东北声望，双方冲突成为区域政治的持续因素。就地方尺度而言，事件并不只属于朝廷或统帅：征发、道路、仓储、户籍、市场、寺院与家庭关系都可能决定命令如何抵达、战事如何延长，或接收何以在不同地区呈现不一样的速度。

本条列举的核心材料为《晋书·慕容皝载记》；《三国史记·高句丽本纪》；《资治通鉴·晋纪》。这些材料较适合钉住事件的发生、官方表述或特定人物、地点，却不自动构成对全境人口、收入、兵数、信仰或动机的调查。账本的置信与材料提示是“远征与丸都受创可靠；路线、俘获人数、王室损失和双方叙事差异有争议”；来源限度进一步说明“前燕—高句丽战争、丸都城与跨境史料互证研究”。因此，不能把条目中的政权名称、行政分类或战报修辞直接转译为固定族群和统一社会意志，也不能将制度宣布写成各地同步实施。

在叙事上，更稳妥的做法是把这一事件视为一条需要地方中介才能运作的链：中央的命令、地方官与精英的选择、运输与劳力的条件、以及普通家庭可能面对的迁徙、赋役或安全风险彼此相连。现有材料能够显示其中若干环节，却保留了大量沉默。承认这种沉默并非放弃解释，而是避免以一条编年记录覆盖不同区域和社群的差异。

账本条目\texttt{JH-0346-HUAN-WEN-INVADES-CHENGHAN}记载，永和二年（346年），东晋与成汉发生“桓温受命自荆州西征成汉，沿长江上游进入巴蜀”。这一节点可放在成汉内乱削弱成都防御，东晋希望以西征提升政权威望并控制上游的条件下理解；其后续影响被账本概括为东晋军队突破峡江防线，巴蜀政权面临终局性军事压力。就地方尺度而言，事件并不只属于朝廷或统帅：征发、道路、仓储、户籍、市场、寺院与家庭关系都可能决定命令如何抵达、战事如何延长，或接收何以在不同地区呈现不一样的速度。

本条列举的核心材料为《晋书·桓温传》；《晋书·穆帝纪》；《华阳国志·李特雄期寿势志》；《资治通鉴·晋纪》。这些材料较适合钉住事件的发生、官方表述或特定人物、地点，却不自动构成对全境人口、收入、兵数、信仰或动机的调查。账本的置信与材料提示是“出师、主要路线和桓温统帅地位可靠；兵数、补给和朝廷授命范围有争议”；来源限度进一步说明“桓温伐蜀、长江上游交通与东晋西部战略研究”。因此，不能把条目中的政权名称、行政分类或战报修辞直接转译为固定族群和统一社会意志，也不能将制度宣布写成各地同步实施。

在叙事上，更稳妥的做法是把这一事件视为一条需要地方中介才能运作的链：中央的命令、地方官与精英的选择、运输与劳力的条件、以及普通家庭可能面对的迁徙、赋役或安全风险彼此相连。现有材料能够显示其中若干环节，却保留了大量沉默。承认这种沉默并非放弃解释，而是避免以一条编年记录覆盖不同区域和社群的差异。

账本条目\texttt{JH-0347-CHENGHAN-FALL}记载，永和三年（347年），东晋与成汉发生“桓温攻入成都，李势投降，成汉灭亡，东晋恢复益州控制”。这一节点可放在成汉军政因继承冲突而弱，桓温军利用长江水陆通道直逼成都的条件下理解；其后续影响被账本概括为东晋重获巴蜀税源与上游战略纵深，桓温的个人威望显著上升。就地方尺度而言，事件并不只属于朝廷或统帅：征发、道路、仓储、户籍、市场、寺院与家庭关系都可能决定命令如何抵达、战事如何延长，或接收何以在不同地区呈现不一样的速度。

本条列举的核心材料为《晋书·穆帝纪》；《晋书·桓温传》；《晋书·李势载记》；《华阳国志·李特雄期寿势志》。这些材料较适合钉住事件的发生、官方表述或特定人物、地点，却不自动构成对全境人口、收入、兵数、信仰或动机的调查。账本的置信与材料提示是“成都降服与成汉灭亡可靠；李势投降条件、战后安置和东晋对蜀地的实际整合有争议”；来源限度进一步说明“成汉灭亡、东晋收复益州与巴蜀战后治理研究”。因此，不能把条目中的政权名称、行政分类或战报修辞直接转译为固定族群和统一社会意志，也不能将制度宣布写成各地同步实施。

在叙事上，更稳妥的做法是把这一事件视为一条需要地方中介才能运作的链：中央的命令、地方官与精英的选择、运输与劳力的条件、以及普通家庭可能面对的迁徙、赋役或安全风险彼此相连。现有材料能够显示其中若干环节，却保留了大量沉默。承认这种沉默并非放弃解释，而是避免以一条编年记录覆盖不同区域和社群的差异。

账本条目\texttt{JH-0349-SHI-HU-DEATH}记载，永宁元年（349年），后赵发生“石虎去世，诸子与养孙冉闵围绕都城邺城和军队展开继承争夺”。这一节点可放在石虎长期以清洗维系统治，未能形成稳定的继承与军权交接机制的条件下理解；其后续影响被账本概括为后赵中央迅速分裂，华北各地军队和地方集团获得自立空间。就地方尺度而言，事件并不只属于朝廷或统帅：征发、道路、仓储、户籍、市场、寺院与家庭关系都可能决定命令如何抵达、战事如何延长，或接收何以在不同地区呈现不一样的速度。

本条列举的核心材料为《晋书·石虎载记》；《晋书·冉闵载记》；《资治通鉴·晋纪》。这些材料较适合钉住事件的发生、官方表述或特定人物、地点，却不自动构成对全境人口、收入、兵数、信仰或动机的调查。账本的置信与材料提示是“石虎卒与继承危机可靠；遗诏、石世即位程序和冉闵早期角色有争议”；来源限度进一步说明“后赵崩解、邺城军队与继承政治研究”。因此，不能把条目中的政权名称、行政分类或战报修辞直接转译为固定族群和统一社会意志，也不能将制度宣布写成各地同步实施。

在叙事上，更稳妥的做法是把这一事件视为一条需要地方中介才能运作的链：中央的命令、地方官与精英的选择、运输与劳力的条件、以及普通家庭可能面对的迁徙、赋役或安全风险彼此相连。现有材料能够显示其中若干环节，却保留了大量沉默。承认这种沉默并非放弃解释，而是避免以一条编年记录覆盖不同区域和社群的差异。

账本条目\texttt{JH-0350-RAN-MIN-WEI}记载，永宁二年（350年），冉魏与后赵遗众发生“冉闵废石祗后在邺城称帝，建立冉魏并在内战中针对羯胡及后赵旧部实施暴力清洗”。这一节点可放在后赵继承战使冉闵掌握邺城军队，他以排斥羯胡和恢复魏号争取部分支持的条件下理解；其后续影响被账本概括为暴力加剧人口逃散和政治碎片化，冉魏同时遭前燕、前秦等多方压力。就地方尺度而言，事件并不只属于朝廷或统帅：征发、道路、仓储、户籍、市场、寺院与家庭关系都可能决定命令如何抵达、战事如何延长，或接收何以在不同地区呈现不一样的速度。

本条列举的核心材料为《晋书·冉闵载记》；《资治通鉴·晋纪》；《十六国春秋》辑佚。这些材料较适合钉住事件的发生、官方表述或特定人物、地点，却不自动构成对全境人口、收入、兵数、信仰或动机的调查。账本的置信与材料提示是“称帝与后赵旧部受大规模杀害可靠；受害群体界定、死亡数字和各地暴力的因果关系有争议”；来源限度进一步说明“冉魏、族群标签、内战暴力与后赵解体研究”。因此，不能把条目中的政权名称、行政分类或战报修辞直接转译为固定族群和统一社会意志，也不能将制度宣布写成各地同步实施。

在叙事上，更稳妥的做法是把这一事件视为一条需要地方中介才能运作的链：中央的命令、地方官与精英的选择、运输与劳力的条件、以及普通家庭可能面对的迁徙、赋役或安全风险彼此相连。现有材料能够显示其中若干环节，却保留了大量沉默。承认这种沉默并非放弃解释，而是避免以一条编年记录覆盖不同区域和社群的差异。

账本条目\texttt{JH-0350-FU-HONG-CLAIMS-QIN}记载，永宁二年（350年），前秦前身发生“氐族首领苻洪在关中东部聚众，自称三秦王，前秦政权的军事基础开始形成”。这一节点可放在后赵崩解使关中东部失去稳定统治，苻洪部可吸纳流民和旧军人的条件下理解；其后续影响被账本概括为苻氏取得建国基础，关中随后成为前秦与其他政权竞争的核心。就地方尺度而言，事件并不只属于朝廷或统帅：征发、道路、仓储、户籍、市场、寺院与家庭关系都可能决定命令如何抵达、战事如何延长，或接收何以在不同地区呈现不一样的速度。

本条列举的核心材料为《晋书·苻洪载记》；《晋书·苻健载记》；《资治通鉴·晋纪》。这些材料较适合钉住事件的发生、官方表述或特定人物、地点，却不自动构成对全境人口、收入、兵数、信仰或动机的调查。账本的置信与材料提示是“苻洪称号与聚众可靠；部众来源、与冉魏和东晋的关系及控制区域有争议”；来源限度进一步说明“前秦建国前史、氐族军团与关中权力真空研究”。因此，不能把条目中的政权名称、行政分类或战报修辞直接转译为固定族群和统一社会意志，也不能将制度宣布写成各地同步实施。

在叙事上，更稳妥的做法是把这一事件视为一条需要地方中介才能运作的链：中央的命令、地方官与精英的选择、运输与劳力的条件、以及普通家庭可能面对的迁徙、赋役或安全风险彼此相连。现有材料能够显示其中若干环节，却保留了大量沉默。承认这种沉默并非放弃解释，而是避免以一条编年记录覆盖不同区域和社群的差异。

账本条目\texttt{JH-0351-FU-JIAN-FOUNDS-FORMER-QIN}记载，永宁三年（351年），前秦发生“苻健在长安称天王，建立前秦并开始重建关中官署”。这一节点可放在苻洪死后苻健继承部众，后赵与冉魏无力稳定关中的条件下理解；其后续影响被账本概括为前秦以长安为基地进入十六国竞争，并逐步吸纳关中人口与官僚。就地方尺度而言，事件并不只属于朝廷或统帅：征发、道路、仓储、户籍、市场、寺院与家庭关系都可能决定命令如何抵达、战事如何延长，或接收何以在不同地区呈现不一样的速度。

本条列举的核心材料为《晋书·苻健载记》；《资治通鉴·晋纪》；《十六国春秋》辑佚。这些材料较适合钉住事件的发生、官方表述或特定人物、地点，却不自动构成对全境人口、收入、兵数、信仰或动机的调查。账本的置信与材料提示是“称天王、入长安和前秦建国可靠；官制规模、城内接收过程与地方归附程度有争议”；来源限度进一步说明“前秦建国、长安重建与关中氐族政权研究”。因此，不能把条目中的政权名称、行政分类或战报修辞直接转译为固定族群和统一社会意志，也不能将制度宣布写成各地同步实施。

在叙事上，更稳妥的做法是把这一事件视为一条需要地方中介才能运作的链：中央的命令、地方官与精英的选择、运输与劳力的条件、以及普通家庭可能面对的迁徙、赋役或安全风险彼此相连。现有材料能够显示其中若干环节，却保留了大量沉默。承认这种沉默并非放弃解释，而是避免以一条编年记录覆盖不同区域和社群的差异。

\subsection{信仰、知识与物质空间}

账本条目\texttt{JH-0352-RAN-MIN-DEFEATED}记载，永和八年（352年），冉魏与前燕发生“冉闵在魏昌一带败于前燕，被俘后处死，冉魏政权迅速瓦解”。这一节点可放在冉魏多面受敌且粮饷不足，前燕利用骑兵和河北联盟发动进攻的条件下理解；其后续影响被账本概括为前燕进入中原政治竞争，后赵崩解后华北并未恢复稳定统治。就地方尺度而言，事件并不只属于朝廷或统帅：征发、道路、仓储、户籍、市场、寺院与家庭关系都可能决定命令如何抵达、战事如何延长，或接收何以在不同地区呈现不一样的速度。

本条列举的核心材料为《晋书·冉闵载记》；《晋书·慕容俊载记》；《资治通鉴·晋纪》。这些材料较适合钉住事件的发生、官方表述或特定人物、地点，却不自动构成对全境人口、收入、兵数、信仰或动机的调查。账本的置信与材料提示是“冉闵被俘处死和冉魏衰亡可靠；战场细节、俘后处置和民间支持范围有争议”；来源限度进一步说明“冉魏覆亡、前燕扩张与华北战后人口流动研究”。因此，不能把条目中的政权名称、行政分类或战报修辞直接转译为固定族群和统一社会意志，也不能将制度宣布写成各地同步实施。

在叙事上，更稳妥的做法是把这一事件视为一条需要地方中介才能运作的链：中央的命令、地方官与精英的选择、运输与劳力的条件、以及普通家庭可能面对的迁徙、赋役或安全风险彼此相连。现有材料能够显示其中若干环节，却保留了大量沉默。承认这种沉默并非放弃解释，而是避免以一条编年记录覆盖不同区域和社群的差异。

账本条目\texttt{JH-0353-LANTING-GATHERING}记载，永和九年三月三日（353年），东晋发生“王羲之与会稽士人在兰亭修禊并作序，成为东晋士人交游与书写文化的标志性事件”。这一节点可放在会稽士族在相对安定的江南组织礼俗与文会，书写成为维系名望网络的媒介的条件下理解；其后续影响被账本概括为兰亭文本在后世反复摹写和政治化，亦提示东晋文化史不能只由宫廷与战事叙述。就地方尺度而言，事件并不只属于朝廷或统帅：征发、道路、仓储、户籍、市场、寺院与家庭关系都可能决定命令如何抵达、战事如何延长，或接收何以在不同地区呈现不一样的速度。

本条列举的核心材料为《晋书·王羲之传》；《世说新语》及刘孝标注；《兰亭序》传世摹本与刻石材料。这些材料较适合钉住事件的发生、官方表述或特定人物、地点，却不自动构成对全境人口、收入、兵数、信仰或动机的调查。账本的置信与材料提示是“修禊和序文形成大体可靠；与会者名单、作品原貌和传世摹本关系有争议”；来源限度进一步说明“兰亭修禊、东晋士族交游与书法文本传播研究”。因此，不能把条目中的政权名称、行政分类或战报修辞直接转译为固定族群和统一社会意志，也不能将制度宣布写成各地同步实施。

在叙事上，更稳妥的做法是把这一事件视为一条需要地方中介才能运作的链：中央的命令、地方官与精英的选择、运输与劳力的条件、以及普通家庭可能面对的迁徙、赋役或安全风险彼此相连。现有材料能够显示其中若干环节，却保留了大量沉默。承认这种沉默并非放弃解释，而是避免以一条编年记录覆盖不同区域和社群的差异。

账本条目\texttt{JH-0354-HUAN-WEN-NORTHERN-EXPEDITION}记载，永和十年（354年），东晋与前秦发生“桓温自荆州北伐前秦，进至关中外围后因粮运困难撤退”。这一节点可放在收复蜀地后桓温声望上升，东晋希望借前秦未稳扩大北方影响的条件下理解；其后续影响被账本概括为补给线限制显示东晋难以长期经营关中，桓温的北伐成为军阀政治资源。就地方尺度而言，事件并不只属于朝廷或统帅：征发、道路、仓储、户籍、市场、寺院与家庭关系都可能决定命令如何抵达、战事如何延长，或接收何以在不同地区呈现不一样的速度。

本条列举的核心材料为《晋书·桓温传》；《晋书·穆帝纪》；《晋书·苻健载记》；《资治通鉴·晋纪》。这些材料较适合钉住事件的发生、官方表述或特定人物、地点，却不自动构成对全境人口、收入、兵数、信仰或动机的调查。账本的置信与材料提示是“北伐、进军关中和撤军可靠；军队数量、长安周边战况与撤退责任有争议”；来源限度进一步说明“桓温第一次北伐、关中补给与东晋军政研究”。因此，不能把条目中的政权名称、行政分类或战报修辞直接转译为固定族群和统一社会意志，也不能将制度宣布写成各地同步实施。

在叙事上，更稳妥的做法是把这一事件视为一条需要地方中介才能运作的链：中央的命令、地方官与精英的选择、运输与劳力的条件、以及普通家庭可能面对的迁徙、赋役或安全风险彼此相连。现有材料能够显示其中若干环节，却保留了大量沉默。承认这种沉默并非放弃解释，而是避免以一条编年记录覆盖不同区域和社群的差异。

账本条目\texttt{JH-0355-LATER-ZHAO-ENDS}记载，永和十一年（355年），后赵残余与冉魏发生“石祗在襄国被杀，后赵最后的政治中心消失”。这一节点可放在石虎死后的继承战与冉魏崛起不断削弱襄国，石祗无法重建有效联盟的条件下理解；其后续影响被账本概括为后赵正式退出竞争，前燕、前秦和各地军镇争夺其遗留人口与城镇。就地方尺度而言，事件并不只属于朝廷或统帅：征发、道路、仓储、户籍、市场、寺院与家庭关系都可能决定命令如何抵达、战事如何延长，或接收何以在不同地区呈现不一样的速度。

本条列举的核心材料为《晋书·石虎载记下》；《晋书·冉闵载记》；《资治通鉴·晋纪》。这些材料较适合钉住事件的发生、官方表述或特定人物、地点，却不自动构成对全境人口、收入、兵数、信仰或动机的调查。账本的置信与材料提示是“石祗被杀与后赵终结大体可靠；行凶者动机、残余军队去向和终结时间界定有争议”；来源限度进一步说明“后赵末期、襄国政权与华北地方化研究”。因此，不能把条目中的政权名称、行政分类或战报修辞直接转译为固定族群和统一社会意志，也不能将制度宣布写成各地同步实施。

在叙事上，更稳妥的做法是把这一事件视为一条需要地方中介才能运作的链：中央的命令、地方官与精英的选择、运输与劳力的条件、以及普通家庭可能面对的迁徙、赋役或安全风险彼此相连。现有材料能够显示其中若干环节，却保留了大量沉默。承认这种沉默并非放弃解释，而是避免以一条编年记录覆盖不同区域和社群的差异。

账本条目\texttt{JH-0356-HUAN-WEN-RECOVERS-LUOYANG}记载，永和十二年（356年），东晋与前燕发生“桓温北伐进入洛阳，修复晋帝陵并短暂恢复东晋在中原的象征性存在”。这一节点可放在后赵瓦解提供北伐窗口，桓温以恢复晋室旧都与陵寝凝聚政治声望的条件下理解；其后续影响被账本概括为东晋未能建立持久的河南统治，但洛阳成为南北争夺的正统象征。就地方尺度而言，事件并不只属于朝廷或统帅：征发、道路、仓储、户籍、市场、寺院与家庭关系都可能决定命令如何抵达、战事如何延长，或接收何以在不同地区呈现不一样的速度。

本条列举的核心材料为《晋书·桓温传》；《晋书·穆帝纪》；《资治通鉴·晋纪》。这些材料较适合钉住事件的发生、官方表述或特定人物、地点，却不自动构成对全境人口、收入、兵数、信仰或动机的调查。账本的置信与材料提示是“入洛阳与修陵大体可靠；控制范围、停留时间和地方响应程度有争议”；来源限度进一步说明“东晋洛阳北伐、帝陵修复与正统象征研究”。因此，不能把条目中的政权名称、行政分类或战报修辞直接转译为固定族群和统一社会意志，也不能将制度宣布写成各地同步实施。

在叙事上，更稳妥的做法是把这一事件视为一条需要地方中介才能运作的链：中央的命令、地方官与精英的选择、运输与劳力的条件、以及普通家庭可能面对的迁徙、赋役或安全风险彼此相连。现有材料能够显示其中若干环节，却保留了大量沉默。承认这种沉默并非放弃解释，而是避免以一条编年记录覆盖不同区域和社群的差异。

账本条目\texttt{JH-0357-FU-JIAN-ACCESSION}记载，永和十三年（357年），前秦发生“苻坚废杀苻生即位，以王猛等整顿前秦朝政”。这一节点可放在苻生统治中的暴力与不安定引起宗室、官僚和军事集团不满的条件下理解；其后续影响被账本概括为苻坚取得皇位后可推行较持续的行政整顿，前秦开始扩张前的重组。就地方尺度而言，事件并不只属于朝廷或统帅：征发、道路、仓储、户籍、市场、寺院与家庭关系都可能决定命令如何抵达、战事如何延长，或接收何以在不同地区呈现不一样的速度。

本条列举的核心材料为《晋书·苻坚载记上》；《晋书·苻生载记》；《资治通鉴·晋纪》。这些材料较适合钉住事件的发生、官方表述或特定人物、地点，却不自动构成对全境人口、收入、兵数、信仰或动机的调查。账本的置信与材料提示是“废立与苻坚即位可靠；苻生罪状、宫廷政变过程和苻坚早期政策有争议”；来源限度进一步说明“苻坚即位、前秦宫廷政变与关中政治整合研究”。因此，不能把条目中的政权名称、行政分类或战报修辞直接转译为固定族群和统一社会意志，也不能将制度宣布写成各地同步实施。

在叙事上，更稳妥的做法是把这一事件视为一条需要地方中介才能运作的链：中央的命令、地方官与精英的选择、运输与劳力的条件、以及普通家庭可能面对的迁徙、赋役或安全风险彼此相连。现有材料能够显示其中若干环节，却保留了大量沉默。承认这种沉默并非放弃解释，而是避免以一条编年记录覆盖不同区域和社群的差异。

账本条目\texttt{JH-0357-WANG-MENG-REFORMS}记载，永和十三年（357年），前秦发生“苻坚任用王猛参与中枢政务，以整饬吏治和抑制豪强为前秦集权奠基”。这一节点可放在苻坚需要可信的文官处理宗室和地方势力，王猛以严明行政获得支持的条件下理解；其后续影响被账本概括为前秦中枢控制增强，为征服前燕和河西提供了财政与官僚条件。就地方尺度而言，事件并不只属于朝廷或统帅：征发、道路、仓储、户籍、市场、寺院与家庭关系都可能决定命令如何抵达、战事如何延长，或接收何以在不同地区呈现不一样的速度。

本条列举的核心材料为《晋书·王猛传》；《晋书·苻坚载记上》；《资治通鉴·晋纪》。这些材料较适合钉住事件的发生、官方表述或特定人物、地点，却不自动构成对全境人口、收入、兵数、信仰或动机的调查。账本的置信与材料提示是“王猛受重用及其整顿政策大体可靠；具体任官次序、法令执行范围和后世美化程度有争议”；来源限度进一步说明“王猛政治、前秦法制与关中地方豪强研究”。因此，不能把条目中的政权名称、行政分类或战报修辞直接转译为固定族群和统一社会意志，也不能将制度宣布写成各地同步实施。

在叙事上，更稳妥的做法是把这一事件视为一条需要地方中介才能运作的链：中央的命令、地方官与精英的选择、运输与劳力的条件、以及普通家庭可能面对的迁徙、赋役或安全风险彼此相连。现有材料能够显示其中若干环节，却保留了大量沉默。承认这种沉默并非放弃解释，而是避免以一条编年记录覆盖不同区域和社群的差异。

账本条目\texttt{JH-0360-MURONG-JUN-EMPEROR}记载，升平四年（360年），前燕发生“慕容儁在邺称帝，前燕由东北政权转为据有中原部分地区的帝国”。这一节点可放在前燕在击败冉魏后控制河北要地，需要以帝号和官僚体系整合新占地区的条件下理解；其后续影响被账本概括为邺城再次成为北方政治中心，前燕与前秦的竞争更趋直接。就地方尺度而言，事件并不只属于朝廷或统帅：征发、道路、仓储、户籍、市场、寺院与家庭关系都可能决定命令如何抵达、战事如何延长，或接收何以在不同地区呈现不一样的速度。

本条列举的核心材料为《晋书·慕容儁载记》；《资治通鉴·晋纪》；《十六国春秋》辑佚。这些材料较适合钉住事件的发生、官方表述或特定人物、地点，却不自动构成对全境人口、收入、兵数、信仰或动机的调查。账本的置信与材料提示是“称帝和以邺为中心可靠；都城营建、官制覆盖和中原士族参与程度有争议”；来源限度进一步说明“前燕帝制、邺城与中原政权转型研究”。因此，不能把条目中的政权名称、行政分类或战报修辞直接转译为固定族群和统一社会意志，也不能将制度宣布写成各地同步实施。

在叙事上，更稳妥的做法是把这一事件视为一条需要地方中介才能运作的链：中央的命令、地方官与精英的选择、运输与劳力的条件、以及普通家庭可能面对的迁徙、赋役或安全风险彼此相连。现有材料能够显示其中若干环节，却保留了大量沉默。承认这种沉默并非放弃解释，而是避免以一条编年记录覆盖不同区域和社群的差异。

账本条目\texttt{JH-0361-AI-EMPEROR-ACCESSION}记载，升平五年（361年），东晋发生“穆帝去世，哀帝司马丕即位，东晋朝政继续由会稽王司马昱等辅政集团主导”。这一节点可放在幼主或病弱皇帝的继承使东晋依赖皇族与高门共同维持中枢运作的条件下理解；其后续影响被账本概括为皇权持续受辅政结构限制，北伐与财政政策难以形成稳定方向。就地方尺度而言，事件并不只属于朝廷或统帅：征发、道路、仓储、户籍、市场、寺院与家庭关系都可能决定命令如何抵达、战事如何延长，或接收何以在不同地区呈现不一样的速度。

本条列举的核心材料为《晋书·穆帝纪》；《晋书·哀帝纪》；《晋书·会稽文孝王昱传》。这些材料较适合钉住事件的发生、官方表述或特定人物、地点，却不自动构成对全境人口、收入、兵数、信仰或动机的调查。账本的置信与材料提示是“穆帝卒与哀帝继位可靠；辅政分工、皇室成员实际权力和士族立场有争议”；来源限度进一步说明“东晋中期继承、辅政集团与门阀政治研究”。因此，不能把条目中的政权名称、行政分类或战报修辞直接转译为固定族群和统一社会意志，也不能将制度宣布写成各地同步实施。

在叙事上，更稳妥的做法是把这一事件视为一条需要地方中介才能运作的链：中央的命令、地方官与精英的选择、运输与劳力的条件、以及普通家庭可能面对的迁徙、赋役或安全风险彼此相连。现有材料能够显示其中若干环节，却保留了大量沉默。承认这种沉默并非放弃解释，而是避免以一条编年记录覆盖不同区域和社群的差异。

账本条目\texttt{JH-0362-MURONG-JUN-DEATH}记载，隆和元年（362年），前燕发生“慕容儁去世，年幼的慕容暐继位，慕容恪等辅政”。这一节点可放在前燕扩张过快而新占地区尚待整合，幼帝继位需要倚重宗室名将的条件下理解；其后续影响被账本概括为慕容恪维持政权一时稳定，但继承弱点为前秦后来进攻留下空间。就地方尺度而言，事件并不只属于朝廷或统帅：征发、道路、仓储、户籍、市场、寺院与家庭关系都可能决定命令如何抵达、战事如何延长，或接收何以在不同地区呈现不一样的速度。

本条列举的核心材料为《晋书·慕容儁载记》；《晋书·慕容暐载记》；《资治通鉴·晋纪》。这些材料较适合钉住事件的发生、官方表述或特定人物、地点，却不自动构成对全境人口、收入、兵数、信仰或动机的调查。账本的置信与材料提示是“慕容儁卒、慕容暐继位和慕容恪辅政可靠；遗诏、辅政权力分配和皇族关系有争议”；来源限度进一步说明“前燕继承、慕容恪辅政与中原北方政局研究”。因此，不能把条目中的政权名称、行政分类或战报修辞直接转译为固定族群和统一社会意志，也不能将制度宣布写成各地同步实施。

在叙事上，更稳妥的做法是把这一事件视为一条需要地方中介才能运作的链：中央的命令、地方官与精英的选择、运输与劳力的条件、以及普通家庭可能面对的迁徙、赋役或安全风险彼此相连。现有材料能够显示其中若干环节，却保留了大量沉默。承认这种沉默并非放弃解释，而是避免以一条编年记录覆盖不同区域和社群的差异。

账本条目\texttt{JH-0365-FORMER-YAN-TAKES-LUOYANG}记载，兴宁三年（365年），前燕与东晋发生“慕容恪攻取洛阳，前燕控制中原旧都并迫使东晋北方据点后撤”。这一节点可放在东晋北伐资源不足且内政不稳，前燕可从河北向河南持续推进的条件下理解；其后续影响被账本概括为洛阳再次脱离东晋控制，南北政权的黄河界线进一步南移。就地方尺度而言，事件并不只属于朝廷或统帅：征发、道路、仓储、户籍、市场、寺院与家庭关系都可能决定命令如何抵达、战事如何延长，或接收何以在不同地区呈现不一样的速度。

本条列举的核心材料为《晋书·慕容恪载记》；《晋书·哀帝纪》；《资治通鉴·晋纪》。这些材料较适合钉住事件的发生、官方表述或特定人物、地点，却不自动构成对全境人口、收入、兵数、信仰或动机的调查。账本的置信与材料提示是“取洛阳与东晋失守可靠；战役路线、守军规模和城市控制范围有争议”；来源限度进一步说明“前燕取洛阳、中原城市与东晋北方战略研究”。因此，不能把条目中的政权名称、行政分类或战报修辞直接转译为固定族群和统一社会意志，也不能将制度宣布写成各地同步实施。

在叙事上，更稳妥的做法是把这一事件视为一条需要地方中介才能运作的链：中央的命令、地方官与精英的选择、运输与劳力的条件、以及普通家庭可能面对的迁徙、赋役或安全风险彼此相连。现有材料能够显示其中若干环节，却保留了大量沉默。承认这种沉默并非放弃解释，而是避免以一条编年记录覆盖不同区域和社群的差异。
